% !TeX root = ../../../main.tex
\subsection{Related Work}

\ChSevenParagraph{Temporal Information Extraction from Clinical Text:}
Extracting timelines from clinical narratives is a challenging biomedical NLP task. The i2b2 2012 challenge introduced annotated datasets for temporal relation extraction from discharge summaries \citep{ch07:uzuner20112010}. Subsequent methods linked clinical events to timestamps or temporal expressions \citep{leeuwenberg2020towards,ch07:frattallone2024using}, typically assuming pre-defined event spans. We adopt our prior approach \citep{ch07:noroozizadeh2025reconstructing,noroozizadeh2025pmoa} to assign timepoints to findings from full-length case reports, enabling finer temporal resolution. Unlike methods using structured EHR data, we focus solely on text—crucial for sources like PubMed that lack structured metadata. By directly supervising event-time alignment, we overcome limitations of span-based annotations~\citep{rosenbloom2011data}. This aligns with growing emphasis on temporality in sepsis and critical care phenotyping \citep{johnson2018comparative,kamran2024evaluation}, particularly given high missingness rates in structured data \citep{ch07:johnson2023mimic,seinen2025using}.

\ChSevenParagraph{Predictive Modeling with Clinical Text:}
Clinical text has been used for outcome prediction tasks like mortality and readmission, with models like \textit{ClinicalBERT} achieving high levels of performance on EHR notes \citep{huang2019clinicalbert,gu2021domain}. 
Text captures complementary information—symptoms and social factors—not found in structured codes. 
However, traditional approaches using full-text or bag-of-words representations obscure temporal dynamics. 
Our approach preserves event order within narratives, treating data as time series. This aligns with recent efforts adapting LLMs for time series tasks, such as Time-LLM mapping numeric sequences to tokens \citep{jin2023time}. 
To our knowledge, this is the first work to forecast from LLM-derived time series in clinical domains. 
While prior sepsis prediction models rely on vitals or scores like SOFA and SAPS \citep{hou2020predicting,noroozizadeh2023temporal}, our method is complementary—leveraging narrative descriptions that include clinician interpretations and context. 
This is valuable when case reports or patient histories are available but structured data is sparse or unavailable.

\ChSevenParagraph{Large Language Models in Healthcare:}
Recent LLMs such as \texttt{GPT-3} and \texttt{GPT-4} have enabled applications from medical QA to note summarization. These models encode substantial medical knowledge and reasoning capacity—\texttt{GPT-4} demonstrates strong performance on board exams and can support patient record abstraction for quality reporting \citep{boussina2024large}. However, their use in clinical forecasting remains largely unexplored.
% Out-of-the-box LLMs may hallucinate and lack domain-specific tuning.
Our benchmark tests LLMs via zero- and few-shot prompting to assess clinical event forecasting from patient narratives, comparing results to fine-tuned models to quantify the gap between general-purpose and task-specific modeling. This contributes evidence that while LLMs hold broad medical knowledge, fine-tuning and structured inputs are often required for clinical reliability \citep{huang2019clinicalbert,gu2021domain}. Some studies suggest domain-specific tuning does not always improve over foundation models \citep{jeong2024medical}, reinforcing the need for task-specific evaluations like ours.


%%%%%%%%%%%%%%%%%%%%%%%%%%%%%%%%%%%%%%%%%%%%%
%%%%%%%%%%%%%%% Methods %%%%%%%%%%%%%%%%
%%%%%%%%%%%%%%%%%%%%%%%%%%%%%%%%%%%%%%%%%%%%%
